\chapter*{生成AI等の利用に関する申告}
\markboth{生成AI等の利用に関する申告}{生成AI等の利用に関する申告}
\addcontentsline{toc}{chapter}{生成AI等の利用に関する申告}

\UsamiFictionalExampleNotice
\par\vspace{1.5\baselineskip}

本論文の作成にあたり，著者はOpenAI ChatGPT（GPT-5，2026年4月1日から7月15日）を，研究課題の論点整理，Pythonコードの実装補助とリファクタリング，章構成の検討，文章表現の推敲および反証観点の洗い出しに使用した．生成AIの出力はそのまま研究上の根拠として採用せず，実験設計，ソースコード，実行結果，数値，図表，主張，引用文献および最終稿を，著者が原資料，一次文献，実行ログおよび再実行結果と照合した．本架空例では，DICOM，NIfTI，患者情報，症例対応表，非公開URLおよび認証情報を外部AIサービスへ入力していない．生成AIを著者として扱わず，本論文の正確性，研究公正，独創性および最終内容については著者が責任を負う．

AI利用の範囲と人間による検証責任を分けて記録する形式については，宇佐美ら\cite{usami2026darkcurrent}の申告例を参考にした．
